%%%%%%%%%%%%
%
% $Autor: Wings $
% $Datum: 2019-03-05 08:03:15Z $
% $Pfad: ListOfParts.tex $
% $Version: 4250 $
% !TeX spellcheck = en_GB/de_DE
% !TeX encoding = utf8
% !TeX root = manual 
% !TeX TXS-program:bibliography = txs:///biber
%
%%%%%%%%%%%%

\chapter{System Configuration}

This chapter provides instructions for configuring the software and hardware settings on your Jetson Nano J1010 to ensure the Vehicle License Plate Detection System operates smoothly.

\section{Configuring the Camera}
To enable real-time license plate detection, configure the connected camera as follows:

\begin{enumerate}
	\item \textbf{Verify Camera Connection:} Ensure the camera is securely connected to the Jetson Nano via the designated camera port.
	\item \textbf{Install Camera Drivers:} If necessary, install drivers for the camera using the command:
	\begin{verbatim}
		sudo apt-get install v4l-utils
	\end{verbatim}
	\item \textbf{Test Camera Functionality:} Use the following command to verify that the camera is functioning:
	\begin{verbatim}
		gst-launch-1.0 nvarguscamerasrc ! nvoverlaysink
	\end{verbatim}
\end{enumerate}

%%%%%%%%%%%%%        PROBABLY TOO MUCH INFORMATION FOR SUCH A MANUAL
%
%\section{Setting Up GPIO for the LED Indicator}
%The LED indicator signals successful license plate detection. To configure the GPIO pins:
%
%\begin{enumerate}
%	\item \textbf{Enable GPIO Access:} Install the GPIO library for Python if it hasn't been installed:
%	\begin{verbatim}
%		sudo pip3 install RPi.GPIO
%	\end{verbatim}
%	\item \textbf{Configure GPIO Pin:} Choose a GPIO pin to connect the LED and configure it in your Python code. For instance, if using pin 18, initialize it as follows:
%	\begin{verbatim}
%		import RPi.GPIO as GPIO
%		GPIO.setmode(GPIO.BOARD)
%		GPIO.setup(18, GPIO.OUT)
%	\end{verbatim}
%\end{enumerate}

\section{Loading the License Plate Detection Model}
Load and configure the pre-trained license plate detection model on the Jetson Nano:

\begin{enumerate}
	\item \textbf{Model Download:} Ensure the trained model file (e.g., `.h5` or `.pt`) is in the project directory.
	\item \textbf{Load Model in Python:} Use TensorFlow or PyTorch to load the model. For TensorFlow:
	\begin{verbatim}
		import tensorflow as tf
		model = tf.keras.models.load_model('path_to_model.h5')
	\end{verbatim}
	For PyTorch:
	\begin{verbatim}
		import torch
		model = torch.load('path_to_model.pt')
		model.eval()
	\end{verbatim}
	\item \textbf{Configure Model Input}: Ensure the model’s input dimensions match those of the camera feed. Resize images if necessary before passing them to the model.
\end{enumerate}

\section{System Test}
Run a basic test to confirm that the system is correctly configured:

\begin{enumerate}
	\item \textbf{Camera Feed Verification:} Verify that the camera feed displays properly and that images can be processed.
	\item \textbf{Model Inference Test:} Run a sample inference on a test image to ensure the model detects the license plate correctly.
	\item \textbf{LED Indicator Test:} Confirm that the LED lights up when a license plate is detected.
\end{enumerate}

After completing these configuration steps, your system should be fully operational. In the next chapter, we will cover operational instructions for using the system.